\ExerciseSection

\begin{enumerate}
\def\labelenumi{\arabic{enumi})}
\item
  设 G 是 \(R^{n}\) 中秩为 p 的离散子群，\((a_{i})_{1\leqslant i\leqslant p}\) 是 G 中 p 个点组成的一个自由系。第 1 号定理 1 的证明表明：G 是由 p 个点 \(d^{-1}a_{i}\) （d 为某个整数）所生成的群的一个子群。试不利用定理 1 证明：存在 G 中 p 个点组成的一个自由系 \(b_{i}=\sum_{j=1}^{p}b_{ij}a_{j}\)，使得若 \(x_{i}=\sum_{j=1}^{p}x_{ij}a_{j}\quad(\mathrm{I}\leqslant i\leqslant p)\) 是 G 中 p 个点组成的任意自由系，就有 \(|\det(x_{ij})|\geqslant|\det(b_{ij})|>0\) 。进而通过证明诸 \(b_{i}\) 生成 G，给出定理 1 的一个不依赖于不变因子理论的证明。（用反证法论证：若 \(z=\sum_{i=1}^{p}z_{i}b_{i}\in G\) 且诸 \(z_{i}\) 中有一个不是整数，则存在点 \(u=\sum_{i=1}^{p}u_{i}b_{i}\) 属于由 z 与诸 \(b_{i}\) 生成的群，使得对某个指标 i 有 \(o<u_{i}<1\) ，由此得出矛盾。）
\item
  设 G 是 \(R^{n}\) 的一个离散子群。若 G 是两个子群 H 与 K 的直和，则由 H 与 K 生成的向量子空间的交仅由 o 组成，因而 G 的秩等于 H 与 K 的秩之和（注意，对 \(R^{n}\) 的每个离散子群 G，由 G 生成的 \(R^{n}\) 的向量子空间典范地同构于 \(G \otimes_{Z} R\) ）。
\item
  设 G 是 \(R^{n}\) 的一个离散子群。要使 G 的子群 H 是 G 的一个直和项，必须且只须 \(H = V \cap G\) ，其中 V 是 \(R^{n}\) 的一个向量子空间（必要性用习题 2；充分性则注意 H 也是 G 与由 H 生成的向量子空间之交）。
\item
  设 H 与 K 是 \(R^{n}\) 的两个离散子群，使得 \(H + K\) 是闭子群。证明 \(H + K\) 此时是离散的，并且 \(r(\mathrm{H}) + r(\mathrm{K}) = r(\mathrm{H} \cap \mathrm{K}) + r(\mathrm{H} + \mathrm{K})\) 。（设 V 是 \(R^{n}\) 中由 \(H \cap K\) 生成的向量子空间；利用习题 3，把 H 分解为 \(V \cap H\) 与一个离散群 \(H_{1}\) 的直和，把 K 分解为 \(V \cap K\) 与一个离散群 \(K_{1}\) 的直和；然后证明和 \(H_{1} + K_{1}\) 是直和，并利用习题 2。）
\item
  设 G、\(G'\) 是 \(R^{n}\) 的两个闭子群，满足 \(G' \subset G\) ，又设 V（相应地 \(V'\) ）是包含在 G（相应地 \(G'\) ）中的最大向量子空间。证明存在一个与 V 互补的向量子空间 W，使得 G 是 V 与离散群 \(K = W \cap G\) 的直和，并且 \(G'\) 是 \(V \cap G'\) 与 \(K \cap G' = W \cap G'\) 的直和。（若 U 是与 \(V'\) 互补的向量子空间，用习题 3 证明离散群 \(U \cap G'\) 是 \(U \cap V \cap G'\) 与一个离散群 \(K'\) 的直和，并取 W 为包含 \(K'\) 的一个向量子空间。）由此推出：
\end{enumerate}

\begin{enumerate}
\def\labelenumi{\alph{enumi})}
\tightlist
\item
  商群 \(G/G'\) 同构于形如
\end{enumerate}

\[
\mathbf {R} ^ {p} \times \mathbf {T} ^ {q} \times \mathbf {Z} ^ {r} \times \mathbf {F},
\]

的乘积群，其中 F 是有限交换群。

\begin{enumerate}
\def\labelenumi{\alph{enumi})}
\setcounter{enumi}{1}
\tightlist
\item
  形如 \(R^{p} \times T^{q} \times Z^{r} \times F\) （F 有限且交换）的群的每个闭子群和每个豪斯多夫商群都具有同样的形式。
\end{enumerate}

\begin{enumerate}
\def\labelenumi{\arabic{enumi})}
\setcounter{enumi}{5}
\tightlist
\item
  设 H 与 K 是 \(R^{n}\) 的两个闭子群，使得 \(H + K\) 是闭子群。
\end{enumerate}

\begin{enumerate}
\def\labelenumi{\alph{enumi})}
\item
  证明：若 V（相应地 W）是包含在 H（相应地 K）中的最大向量子空间，则 \(V + W\) 是包含在 G 中的最大向量子空间，因而 \(d(\mathrm{H}) + d(\mathrm{K}) = d(\mathrm{H} \cap \mathrm{K}) + d(\mathrm{H} + \mathrm{K})\) {[}注意 \(\mathrm{G}/(\mathrm{V} + \mathrm{W})\) 的有理秩有限，因而 G 不能包含一条不含于 \(V + W\) 中的直线{]}。
\item
  设 U 是与 \(V + W\) 互补的一个子空间，使得 G 是 \(V + W\) 与 \(M = G \cap U\) 的直和，并且 \(H \cap K\) 是 \((V + W) \cap (H \cap K)\) 与 \(L = H \cap K \cap U\) 的直和（习题 5）。设 \(H'\) （相应地 \(K'\) ）是 M 中由 H（相应地 K）的点在 G 分解为 \(V + W\) 与 M 的直和中的分量所组成的子群。证明 \(M = H' + K'\) ，并且 \(H' \cap K' = L\) 。
\item
  若令 \(H'' = H \cap (V + W)\) ，\(K'' = K \cap (V + W)\) ，证明 \(r(H'') + r(K'') = r(H'' \cap K'') + r(H'' + K'')\) （化归为 \(V \cap W = \{o\}\) 的情形，并证明此时 \(H'' \cap K''\) 是 \(V \cap K''\) 与 \(W \cap H''\) 的直和。
\end{enumerate}

\(d\) ) 证明

\[
r (\mathrm{H}) + r (\mathrm{K}) = r (\mathrm{H} \cap \mathrm{K}) + r (\mathrm{H} + \mathrm{K})
\]

{[}利用 c) 与习题 4，注意 \(r(\mathbf{H}) = r(\mathbf{H}') + r(\mathbf{H}'')\) 以及 \(r(\mathbf{K}) = r(\mathbf{K}') + r(\mathbf{K}'' )\) {]}。

\begin{enumerate}
\def\labelenumi{\arabic{enumi})}
\setcounter{enumi}{6}
\item
  设 G 是 \(R^{n}\) 的一个闭子群，H 是 G 的一个闭子群。则 G 是 H 与另一个闭子群 K 的直和，当且仅当 H 是 G 与某个向量子空间之交。（必要性用习题 5；为证充分性，对 G 与 H 适当地应用第 2 号的定理 2，并利用习题 3）。
\item
  \begin{enumerate}
  \def\labelenumii{\alph{enumii})}
  \tightlist
  \item
    设 H、K 是 \(R^{n}\) 的两个闭子群。证明：若 \(r(\mathrm{H}) + r(\mathrm{K})\) 等于 \(r(\mathrm{H} \cap \mathrm{K}) + r(\mathrm{H} + \mathrm{K})\) ，则 \(H + K\) 在 \(R^{n}\) 中是闭的。{[}利用习题 7 与所给假设，化归为 H、K 与 \(H \cap K\) 都生成同一向量子空间的情形；若 V（相应地 W）表示包含在 H（相应地 K）中的最大子空间，则借助习题 5 证明 V 由 \(V \cap K\) 生成，W 由 \(W \cap H\) 生成；最后利用习题 7，把 \(H \cap K\) 分解为 \(H \cap K \cap (V + W)\) 与一个离散群的直和。{]}
  \end{enumerate}
\end{enumerate}

\begin{enumerate}
\def\labelenumi{\alph{enumi})}
\setcounter{enumi}{1}
\item
  由 a) 与习题 6 推出：若 H 与 K 是 \(R^{n}\) 的两个闭子群，使得 \(H + K\) 是闭子群，则 \(H^{*} + K^{*}\) 也是闭子群。
\item
  由 b) 推出：若 H 与 K 是 \(R^{n}\) 的两个闭子群，使得 \(d(\mathrm{H}) + d(\mathrm{K}) = d(\mathrm{H} \cap \mathrm{K}) + d(\overline{\mathrm{H} + \mathrm{K}})\) ，则子群 \(H + K\) 是闭的。
\end{enumerate}

\begin{enumerate}
\def\labelenumi{\arabic{enumi})}
\setcounter{enumi}{8}
\item
  设 G 是 \(R^{n}\) 的一个子群，不必是闭的，秩为 p，又设 V 是包含在 \(\overline{G}\) 中的最大向量子空间。若 V 的维数 \(q \leqslant p\) ，证明 G 是 \(V \cap G\) 与一个秩为 p - q 的离散子群的直和，该离散子群包含在与 V 互补的一个向量子空间中{[}注意，对每个 \(x \in \overline{G}\) ，\((x + V) \cap G\) 在线性簇 \(x + V\) 中稠密{]}。
\item
  设 \(a\) 是一个实数，\(n\) 是一个 \(>0\) 的整数，考虑数 \(x_{k} = ka - [ka] \in [0, \mathrm{I}[(\mathrm{I} \leqslant k \leqslant n + \mathrm{I})\) 。若 \(\mathbf{I}_h\) 表示区间
\end{enumerate}

\[
\left[ \frac {h - \mathrm{I}}{n}, \frac {h}{n} \right] \quad (\mathrm{I} \leqslant h \leqslant n),
\]

证明存在两个不同的指标 \(k, k'\) ，使得 \(x_k\) 与 \(x_{k'}\) 属于同一个区间 \(I_h\) （对某个适当的 \(h\) 值）。由此推出存在两个整数 \(p, q\) ，使得 \(1 \leqslant p \leqslant n\) 且 \(|pa - q| \leqslant 1 / n\) 。

\begin{enumerate}
\def\labelenumi{\Roman{enumi})}
\setcounter{enumi}{1}
\tightlist
\item
  设 \(\boldsymbol{a}_i = (a_{ij})\) （ \(1 \leqslant i \leqslant m\) ，\(1 \leqslant j \leqslant n\) ）是 \(m\) 个属于 \(\mathbf{R}^n\) 的点，又设 \(q\) 是一个 \(>0\) 的整数。对每个形如 \(\boldsymbol{k} = (k_i)\) （ \(1 \leqslant i \leqslant m\) ）的点，它属于 \(\mathbf{R}^m\) 、具有整数坐标（不全为零）且满足不等式 \(0 \leqslant k_i \leqslant q\) ；对这样的点，考虑点 \(x_k = (x_{j,k})\) （ \(1 \leqslant j \leqslant n\) ），它属于 \(\mathbf{R}^n\) ，其所有坐标都属于 {[}o, i{[} ，并且它模 \(Z^n\) 同余于 \(\sum_{i=1}^{m} k_i a_i\) ，即坐标为
\end{enumerate}

\[
x _ {j, k} = \sum_ {i = 1} ^ {m} k _ {i} a _ {i j} - \left[ \sum_ {i = 1} ^ {m} k _ {i} a _ {i j} \right].
\]

  的点。若 \(p\) 是使 \(q + 1 \geqslant p^{n/m}\) 的最小整数，证明存在两个不同的点 \(k, k'\) ，使得 \(x_k\) 与 \(x_{k'}\) 都属于同一个棱长为 \(1/p\) 的立方体（方法与习题 10 相同）。由此推出：存在 \(m\) 个不全为零的整数 \(p_i\) 以及 \(n\) 个整数 \(r_j\) （ \(1 \leqslant j \leqslant n\) ），使得 \(0 \leqslant p_i \leqslant q\) 对 \(1 \leqslant i \leqslant m\) 成立，并且

\[
\left| \sum_ {i = 1} ^ {m} p _ {i} a _ {i j} - r _ {j} \right| \leqslant \frac {\mathrm{I}}{p} \quad (\mathrm{I} \leqslant j \leqslant n).
\]

利用这一结果，给出第 1 号命题 2（“抽屉原理”）的另一个证明。

\begin{enumerate}
\def\labelenumi{\arabic{enumi})}
\setcounter{enumi}{11}
\tightlist
\item
  对每一对实数 \((\theta, \beta)\) ，存在无穷多个整数三元组 \((p, q, r)\) ，使得 \(r > 0\) ，\(|q| \leqslant \frac{1}{2} r\) ，并且
\end{enumerate}

\[
- \frac {\mathrm{I}}{r} <   \theta q + p - \beta <   \frac {\mathrm{I}}{r}.
\]

{[}若 \(m, n\) 是使 \(|n\theta - m| < 1/n\) 的两个整数（习题 10），取 \(r = n\) ，并选取 \(p, q\) 使得 \(pn + qm\) 与 \(n\beta\) 之差小于 \(1/2\) 。{]}

\begin{enumerate}
\def\labelenumi{\arabic{enumi})}
\setcounter{enumi}{12}
\tightlist
\item
  \(n\) 阶法里数列（\(n\) 为 \(>0\) 的整数）是集 \(F_n\) ，它由既约形式的有理数 \(p/q\) （满足 \(0 \leqslant p \leqslant q \leqslant n\) ）组成，按递增次序排列（也就是说，法里数列是一个序列）。
\end{enumerate}

\begin{enumerate}
\def\labelenumi{\alph{enumi})}
\item
  证明：若两个有理数 \(r = p / q\) 、\(r' = p' / q'\) 满足 \(p'q - pq' = \pm 1\) ，则每一对整数 \((p'', q'')\) 都可以表示成形式 \(p'' = px + p'y\) ，\(q'' = qx + q'y\) ，其中 \(x\) 与 \(y\) 是整数。分数 \(p'' / q''\) 属于以 \(r\) 与 \(r'\) 为端点的闭区间，当且仅当 \(x\) 与 \(y\) 同号。
\item
  由 a) 推出：若 \(r = p / q\) 与 \(r' = p' / q'\) 是 {[}0, 1{]} 中的两个有理数，满足 \(q > 0\) 、\(q' > 0\) 且 \(qp' - pq' = \pm 1\) ，则 \(r\) 与 \(r'\) 是数列 \(\mathbf{F}_n\) 中相邻的元素，其中 \(n = \max (q, q')\) 。此外，设 \(m\) 为使以 \(r\) 与 \(r'\) 为端点的开区间含有 \(\mathbf{F}_m\) 中一点的最小整数，则 \(m = q + q'\) ，并且在此区间中只有 \(\mathbf{F}_{q + q'}\) 的一个点，即分数 \((p + p') / (q + q')\) 。
\item
  反之，证明若 \(r\) 与 \(r'\) 是 \(\mathbf{F}_n\) 的两个相邻元素，则有 \(p'q - pq' = \pm 1\) （对 \(n\) 用归纳法）。
\item
  由 c) 推出：对每个实数 \(\theta \in [\mathrm{o},\mathrm{i}]\) 以及每个整数 \(n\geqslant \mathrm{i}\) ，至少存在一个既约形式的分数 \(p / q\) ，使得 \(\mathrm{i}\leqslant q\leqslant n\) 且 \(|\theta -p / q|\leqslant \mathrm{i} / (n + \mathrm{i})q\) （参看习题 10）。
\item
  若 \(p / q\) 是一个既约形式的分数，满足 \(|\theta - p / q| < 1 / q^2\) ，则 \(\theta\) 属于以法里数列 \(\mathbf{F}_q\) 中与 \(p / q\) 相邻的两项为端点的开区间。
\end{enumerate}

\begin{enumerate}
\def\labelenumi{\arabic{enumi})}
\setcounter{enumi}{13}
\tightlist
\item
  设 \(\varphi: \mathbf{R} \to \mathbf{T}\) 是典范同态；设 \(\theta\) 是 \(\mathbf{T}\) 中一个无限阶的元素，又设 \(\theta_0\) 是使 \(\varphi(\theta_0) = \theta\) 的（无理）实数。对每个整数 \(n > 0\) ，设 \(\mathbf{S}_n\) 是形如 \(k\theta\) 的元素（属于 \(\mathbf{T}\) ，且 \(\mathbf{i} \leqslant k \leqslant n\) ）所成的集，又对每个区间 \(\mathbf{I}\) （\(\mathbf{R}\) 的），设 \(\mathrm{N}(\mathbf{I}, n)\) 是集 \(\mathbf{I} \cap \overline{\varphi}^1(\mathbf{S}_n)\) 的元素个数。
\end{enumerate}

\begin{enumerate}
\def\labelenumi{\alph{enumi})}
\item
  若取 \(\mathbf{I} = [\mathrm{o},\theta_0[\) ，证明 \(\mathrm{N(I,n)} = [n\theta_0]\) {[}注意，此时 \(\mathrm{N(I,n)}\) 是如下整数对 \((x,y)\) 的个数：\(\mathrm{i}\leqslant y\leqslant n\) 且 \(x\leqslant y\theta_0 < x + \theta_0]\) 。
\item
  更一般地，若取 \(\mathbf{I} = [\mathrm{o}, m\theta_0[, m\) 为整数，则有 \(m. [(n - m)\theta_0] \leqslant \mathrm{N}(\mathrm{I}, n) \leqslant m. [n\theta_0]\) （同样的方法）。
\item
  由此推出：若 I 是任意区间，则数 \(\mathrm{N}(\mathrm{I}, n)/n\) 当 \(n \to +\infty\) 时趋于一个等于 I 的长度的极限。（先对形如 \([m\theta_{0} + a, m'\theta_{0} + a']\) 的 I 证明该结果，其中 \(m, m', a\) 与 \(a'\) 是整数，利用 b)；然后用形如 \(m\theta_{0} + a\) 的数逼近 I 的端点，过渡到一般情形。）{[}“序列 \((k\theta) \mod 1\) 的均分”。{]}
\end{enumerate}

* 15) 设 I 是 \(\mathbf{R}^n\) 中一个棱长为 \(2\pi\) 的闭立方体。把 I 的每个点 \(x = (x_i)\) 映为点 \(y = (y_j)\) （属于 \(\mathbf{R}^{n+1}\) ），使得

\[
\begin{array}{l} y _ {1} = \sin x _ {1} \\ y _ {2} = (2 + \cos x _ {1}) \sin x _ {2} \end{array}
\]

\[
y _ {p} = \left(2 ^ {p - 1} + \sum_ {k = 1} ^ {p - 1} 2 ^ {p - k - 1} \cos x _ {p - k} \cos x _ {p - k + 1} \dots \cos x _ {p - 1}\right) \sin x _ {p} \quad (2 \leqslant p \leqslant n)
\]

\[
y _ {n + 1} = \sum_ {k = 1} ^ {n - 1} 2 ^ {p - k - 1} \cos x _ {n - k} \cos x _ {n - k + 1} \dots \cos x _ {n}.
\]

 证明 I 在此映射下的像同胚于 \(T^{n}\) （对 n 用归纳法证明：注意 \(y_{p}\) 与 \(\sin x_{p}\) 当 \(2 \leqslant p \leqslant n\) 时同号）。若 n=2，这样定义的 \(R^{3}\) 的子集称为旋转环面。*

¶ 16) 设 G 是 \(\mathbf{R}^n\) 的一个子群，它包含一个紧连通子集 K，使得由 K 生成的仿射线性簇的维数是 \(p\) 。证明 G 包含一个维数为 \(p\) 的向量子空间。{[}对 \(n\) 用归纳法证明，并利用第六章 § 1 的习题 10 c)。{]}

\(\mathbb{J}\) 17) 设 \(\mathbf{E}\) 是乘积 \((\mathbf{Q}_p)^n\) ：它由 \(n\) 个都等于域 \(\mathbf{Q}_p\) （\(p\) 进数域）的因子组成（第三章，§ 6，习题 23），赋有由诸因子 \(\mathbf{Q}_p\) 的加法拓扑结构的乘积所成的拓扑群结构，以及它的 \(n\) 维向量空间结构（在域 \(\mathbf{Q}_p\) 上）。若 \(v\) 表示加法 \(p\) 进赋值（在 \(\mathbf{Q}_p\) 上），则令 \(|x|_p = p^{-v_p(x)}\) ，其中 \(x \neq 0\) 属于 \(\mathbf{Q}_p\) ，并令 \(|o|_p = 0\) ；又令 \(||x|| = \sup |x_i|_p\) （对一切 \(x = (x_i)_{1 \leqslant i \leqslant n} \in \mathbf{E}\) ）（参看第九章，§ 3，第 2 号与第 3 号）。

\begin{enumerate}
\def\labelenumi{\alph{enumi})}
\tightlist
\item
  子集 \(\mathbf{A}\) （\(\mathbf{E}\) 的）是相对紧的，当且仅当
\end{enumerate}

\[
\sup _ {\mathbf {x} \in \mathbf {A}} | | \mathbf {x} | | <   + \infty .
\]

{[}注意 \(v_{p}\) 是 \(\mathbf{Q}_p\) 到 \(\mathbf{Z}\) （带离散拓扑）中的连续映射。{]}

\begin{enumerate}
\def\labelenumi{\alph{enumi})}
\setcounter{enumi}{1}
\item
  若 \(\mathbf{G}\) 是 \(\mathbf{E}\) 的一个闭子群，则 \(\mathbf{G}\) 是环 \(\mathbf{Z}_p\) （\(p\) 进整数环）上的一个拓扑模（若 \(\boldsymbol{x} \in \mathbf{G}\) ，则有 \(n\boldsymbol{x} \in \mathbf{G}\) （对一切 \(n \in \mathbf{Z}\) ），而 \(\mathbf{Z}\) 在 \(\mathbf{Z}_p\) 中稠密）。
\item
  若 K 是 E 的一个紧子群，则存在 E 的 \(m \leqslant n\) 个点 \(a_{i}\) （ \(i \leqslant i \leqslant m\) ）组成的一个自由系，使得 K 是 m 个群 \(Z_{p}.a_{i}\) 的直和。{[}利用 a) 证明：若 \((\boldsymbol{e}_{j})_{1 \leqslant j \leqslant n}\) 是 E 的标准基，则存在一个整数 \(k \in Z\) ，使得 K 包含在 n 个群 \(Z_{p}.p^{k}\boldsymbol{e}_{j}\) 的直和中；然后应用主理想环上的模的理论。{]}
\item
  若 G 是 E 的一个闭（但非紧）子群，则 G 包含一个一维向量子空间，即 \(Q_{p} \cdot a\) ，其中 \(a \neq 0\) 。{[}设 C 是由所有点 \(x \in E\) （满足 \(\|x\| = 1\) ）组成的 E 的紧子集；证明 \(G \cap C\) 包含一个点列 \((\boldsymbol{x}_{r})_{r \in \mathbb{N}}\) ，使得 \(p^{-r}x_{r} \in G\) ，并取 a 为该序列的一个聚点。{]}
\item
  设 G 是 E 的一个闭子群，V 是包含在 G 中的最大向量子空间，W 是与 V 互补的一个向量子空间；证明 \(W \cap G\) 是紧的，并且 G 是 V 与 \(W \cap G\) 的直和（按第 2 号定理 2 的方式论证）。
\end{enumerate}

\(f)\) 给定一个子群 \(\mathbf{G}\) （\(\mathbf{E}\) 的），用 \(\mathbf{G}^*\) 表示所有 \(\boldsymbol{u} = (u_i) \in \mathbf{E}\) 中使得

\[
(u | x) = \sum_ {i = 1} ^ {n} u _ {i} x _ {i}
\]

属于 \(\mathbf{Z}_p\) （对所有 \(x = (x_i) \in \mathbf{G}\) ）者所成的集。证明 \(\mathbf{G}^*\) 在 \(\mathbf{E}\) 中是闭的，\((\overline{\mathbf{G}})^* = \mathbf{G}^*\) ，并且 \((\mathbf{G}^*)^* = \overline{\mathbf{G}}\) 。

\begin{enumerate}
\def\labelenumi{\alph{enumi})}
\setcounter{enumi}{6}
\tightlist
\item
  对 E 的闭子群，叙述并证明习题 2 至 8（包括 8 在内）的类似命题。特别地，形如
\end{enumerate}

\[
(\mathbf {Q} _ {p}) ^ {r} \times (\mathbf {Q} _ {p} / \mathbf {Z} _ {p}) ^ {s} \times (\mathbf {Z} _ {p}) ^ {t} \times \mathbf {F}
\]

（其中 F 是有限个 p 幂阶循环群的乘积）的群的每个闭子群和每个豪斯多夫商群都是同样形式的群。

